% PDF p. 356
\AppendixSection{6J}{反向时间随机微分方程}

\AppendixHeading{6J1}{设定和记号}

设
\begin{equation*}
  dX_t=g(X_t,t)\,dt+\sigma(t)\,dW_t,\qquad t\in[0,T],\quad T=1,
  \tag{6J1}\label{eq:6J1}
\end{equation*}
其中 $g\colon\R^d\times[0,T]\to\R^d$ 关于 $x$ 为 $C^1$，且关于 $(x,t)$ 有界。$\sigma\colon[0,T]\to(0,\infty)$ 为 $C^1$。为简明起见取标量形式，完整的矩阵形式只需调整记号。$W_t$ 是带滤过的概率空间 $(\Omega,\mathcal F,(\mathcal F_t)_{t\geq0},P)$ 上的标准 $d$ 维维纳过程。

扩散过程的反向过程本身也是扩散过程，具有如下反向时间 SDE，
\begin{equation*}
  \begin{aligned}
   dx_t&=\bigl[g(x,t)-\sigma^2(t)\nabla_{x_t}\log f_t(x)\bigr]\,dt
           +\sigma(t)\,d\bar W_t,\\
   x_0&=x(0),\qquad x_1=x(1),
  \end{aligned}
  \tag{6J2}\label{eq:6J2}
\end{equation*}
其中 $\bar W_t$ 为从 $t$ 到 0 沿时间反向运行的标准维纳过程，$dt$ 为负的无穷小时间步长。

假定 $X_t$ 的前向分布有严格为正的 $C^{2,1}$ 密度 $f_t(x)$，满足福克–普朗克方程（前向柯尔莫哥洛夫方程，forward Kolmogorov equation），
\begin{equation*}
  \frac{\partial f_t(x)}{\partial t}
   =-\nabla\cdot\bigl(g(x,t)f_t(x)\bigr)
     +\frac{\sigma^2(t)}2\Delta f_t(x),\qquad f_0(x)=f_X(x).
  \tag{6J3}\label{eq:6J3}
\end{equation*}

\AppendixHeading{6J2}{反向时间坐标}

定义反向时钟 $s=T-t$ 和时间反演过程
\[
  Y_s:=X_{T-s},\qquad s\in[0,T].
\]
由于 $X$ 为马尔可夫过程，$(Y_s)_{s\in[0,T]}$ 关于反向滤过
\[
  \mathcal G_s:=\sigma(Y_r,\ 0\leq r\leq s)
             =\sigma(X_{T-r},\ 0\leq r\leq s)
\]
也是马尔可夫过程。

% PDF p. 357
我们的目标是证明 $Y$ 满足反向时间 SDE，
\begin{equation*}
  dY_s=\bigl(g(Y_s,T-s)-\sigma^2(T-s)\nabla\log f_{T-s}(Y_s)\bigr)\,ds
        +\sigma(T-s)\,d\tilde W_s,
  \tag{6J4}\label{eq:6J4}
\end{equation*}
其中 $\tilde W_s$ 是关于 $\mathcal G_s$ 的标准维纳过程（沿时间反向运行的布朗运动）。

由于 $ds=-dt$，式 \eqref{eq:6J4} 的漂移恰好是式 \eqref{eq:6J2} 中的表达式，并满足“$ds>0$ 对应 $dt<0$”的符号约定。\translatornote{原书式 \eqref{eq:6J4} 沿用了式 \eqref{eq:6J2} 在负时间步 $dt$ 下的漂移系数。改用递增时间 $s$ 后，$ds=-dt$，漂移应为 $-g(Y_s,T-s)+\sigma^2(T-s)\nabla\log f_{T-s}(Y_s)$。}

\AppendixHeading{6J3}{生成元和条件漂移}

\AppendixSubheading{6J3.1}{前向生成元}

对任意 $\varphi\in C_b^2(\R^d)$，定义生成元（generator）
\begin{equation*}
  \mathcal L_t\varphi=g(x,t)\nabla\varphi(x)+\frac{\sigma^2(t)}2\Delta\varphi.
  \tag{6J5}\label{eq:6J5}
\end{equation*}
由伊藤公式（Itô's formula），
\begin{equation*}
  M_t^\varphi=\varphi(X_t)-\varphi(X_0)-\int_0^t\mathcal L_u\varphi(X_u)\,du
  \tag{6J6}\label{eq:6J6}
\end{equation*}
是关于 $(\mathcal F_t)$ 的鞅（martingale）。

\AppendixSubheading{6J3.2}{反向条件期望}

固定 $s\in(0,T)$，令 $\psi\in C_b^2$。由于 $Y_s=X_{T-s}$，有
\begin{equation*}
  E[\psi(Y_{s+h})-\psi(Y_s)\mid\mathcal G_s]
    =E[\psi(X_{T-s-h})-\psi(X_{T-s})\mid X_{T-s}].
  \tag{6J7}\label{eq:6J7}
\end{equation*}
记 $t:=T-s$，取小的 $h>0$。以 $X_t=x$ 为条件，有
\begin{equation*}
  \frac1h E_x[\psi(X_{t-h})-\psi(X_t)]
    =\frac1h\bigl(E_x[\psi(X_{t-h})]-\psi(x)\bigr).
  \tag{6J8}\label{eq:6J8}
\end{equation*}
根据后向柯尔莫哥洛夫方程（backward Kolmogorov equation），当 $h\to0^+$ 时右端为 $-\mathcal L_t\psi(x)+o(1)$。不过，$\mathcal L_t$ 作用于前向转移密度，我们需要一个反向作用的生成元。

设 $f_t(x)$ 是 $X(t)$ 的密度函数。回顾福克–普朗克方程（前向柯尔莫哥洛夫方程），
\begin{equation*}
  \partial_t f_t(x)=-\nabla\cdot\bigl(g(x,t)f_t(x)\bigr)
      +\tfrac12\sigma^2(t)\Delta f_t(x),\qquad f_0(x)=f_X(x).
  \tag{6J9}\label{eq:6J9}
\end{equation*}
将福克–普朗克方程乘以 $\psi(x)$，再分部积分，得到
% PDF p. 358
\begin{equation*}
  \begin{aligned}
   \int_{\R^d}\partial_t f_t(x)\psi(x)\,dx
     &=-\int_{\R^d}\bigl(\nabla\cdot(g(x,t)f_t(x))\bigr)\psi(x)\,dx\\
     &\quad+\int_{\R^d}\tfrac12\sigma^2(t)\Delta f_t(x)\psi(x)\,dx\\
     &=0+\int_{\R^d}\bigl(\nabla\psi(x)\bigr)\cdot(g(x,t)f_t(x))\,dx\\
     &\quad+0-\tfrac12\sigma^2(t)
             \int_{\R^d}\bigl(\nabla\psi(x)\bigr)\cdot\bigl(\nabla f_t(x)\bigr)x\\
     &=\int_{\R^d}\bigl(\nabla\psi(x)\bigr)\cdot(g(x,t)f_t(x))\,dx\\
     &\quad+\tfrac12\sigma^2(t)\int_{\R^d}\Delta\psi(x)f_t(x)\,dx\\
     &=\int_{\R^d}f_t(x)
           \bigl[g(x,t)\cdot\nabla\psi(x)+\tfrac12\sigma^2(t)\Delta\psi(x)\bigr]\,dx.
  \end{aligned}
  \tag{6J10}\label{eq:6J10}
\end{equation*}
\translatornote{式 \eqref{eq:6J10} 第四行末尾的 $x$ 按分部积分应为 $dx$。}
由于 $E[\psi(x)]=\int_{\R^d}\psi(x)f_t(x)\,dx$，式 \eqref{eq:6J10} 可改写为
\begin{equation*}
  \partial_t E[\psi(X_t)]
   =E\bigl[g(x,t)\cdot(\nabla\psi(x))+\tfrac12\sigma^2(t)\Delta\psi(x)\bigr].
  \tag{6J11}\label{eq:6J11}
\end{equation*}
将式 \eqref{eq:6J11} 从 $t-h$ 积分到 $t$，其中 $h>0$ 较小，
\begin{equation*}
  E[\psi(X_t)]-E[\psi(X_{t-h})]
    =h E\bigl[g(x,t)\cdot(\nabla\psi(X_t))+\tfrac12\sigma^2(t)\Delta\psi(X_t)\bigr].
  \tag{6J12}\label{eq:6J12}
\end{equation*}

\AppendixSubheading{6J3.3}{引入条件期望（全期望公式）}

记 $\varphi_h(x)=E[\psi(X_{t-h})\mid X_t=x]$。用条件期望计算，得到
\[
  E[\psi(X_{t-h})]=\int\varphi_h(x)f_t(x)\,dx,\qquad
  E[\psi(X_t)]=\int\psi(x)f_t(x)\,dx.
\]
于是式 \eqref{eq:6J12} 成为
\begin{equation*}
  \begin{aligned}
   \int[\psi(x)-\varphi_h(x)]f_t(x)\,dx
    &=h\int\bigl[g(x,t)\cdot(\nabla\psi(x))
                +\tfrac12\sigma^2(t)\Delta\psi(x)\bigr]f_t(x)\,dx+o(h).
  \end{aligned}
  \tag{6J13}\label{eq:6J13}
\end{equation*}
令
\[
  I=\frac{\sigma^2(t)}2\int f_t(x)\Delta\psi(x)\,dx.
\]
% PDF p. 359
分部积分得到
\begin{equation*}
  \begin{aligned}
    I&=\frac{\sigma^2(t)}2\int f_t(x)\Delta\psi(x)\,dx\\
     &=\frac{\sigma^2(t)}2\int\sum_{i=1}^d f_t\,\partial_{x_i x_i}\psi(x)\,dx\\
     &=-\frac{\sigma^2(t)}2\sum_{i=1}^d\int(\partial_{x_i}f_t)\partial_{x_i}\psi(x)\,dx\\
     &=-\frac{\sigma^2(t)}2\int(\nabla f_t)\cdot\nabla\psi(x)\,gx\\
     &=-\sigma^2(t)\int(\nabla f_t)\cdot\nabla\psi(x)\,gx
           +\frac{\sigma^2(t)}2\int f_t(x)\Delta\psi(x)\,dx\\
     &=-\int f_t(x)\bigl[\sigma^2(t)\nabla\log f_t(x)\bigr]\cdot\nabla\psi(x)\,dx
           +\frac{\sigma^2(t)}2\int f_t(x)\Delta\psi(x)\,dx .
  \end{aligned}
  \tag{6J14}\label{eq:6J14}
\end{equation*}
\translatornote{式 \eqref{eq:6J14} 的两处 $gx$ 按积分记号应为 $dx$。此外，倒数两行把第一行中的 $I$ 加到了分部积分所得的 $2I$ 上，符号和系数不能由前一行得到。}
由式 \eqref{eq:6J13} 回顾
\begin{equation*}
  \begin{aligned}
   \int[\psi(x)-\varphi_h(x)]f_t(x)\,dx
     &=h\int\bigl[g(x,t)\cdot(\nabla\psi(x))
                   +\tfrac12\sigma^2(t)\Delta\psi(x)\bigr]f_t(x)\,dx\\
     &=h\int g(x,t)\cdot(\nabla\psi(x))f_t(x)\,dx+Ih.
  \end{aligned}
  \tag{6J15}\label{eq:6J15}
\end{equation*}
\begin{equation*}
  \begin{aligned}
   &=h\int g(x,t)\cdot(\nabla\psi(x))f_t(x)\,dx
        -h\int f_t(x)\bigl[\sigma^2(t)\nabla\log f_t(x)\bigr]\cdot\nabla\psi(x)\,dx\\
   &\quad+h\int f_t(x)\frac{\sigma^2(t)}2\Delta\psi(x)\,dx.
  \end{aligned}
  \tag{6J16}\label{eq:6J16}
\end{equation*}
由于式 \eqref{eq:6J16} 对每个点都成立，得到
\begin{equation*}
  \begin{aligned}
  \psi(x)-\varphi_h(x)
     &=h\bigl[(g(x,t)-\sigma^2(t)\nabla\log f_t(x))\cdot\nabla\psi(x)\bigr]\\
     &\quad+h\frac{\sigma^2(t)}2\Delta\psi(x).
  \end{aligned}
  \tag{6J17}\label{eq:6J17}
\end{equation*}
\translatornote{上一步给出的是关于 $x$ 的积分等式。仅凭这一等式无法推出逐点的条件漂移，原书推导在此缺少从联合或条件分布出发的论证。}
所以
\begin{equation*}
  \frac{\psi(x)-\varphi_h(x)}{h}
     =\bigl[(g(x,t)-\sigma^2(t)\nabla\log f_t(x))\cdot\nabla\psi(x)\bigr]
      +\frac{\sigma^2(t)}2\Delta\psi(x).
  \tag{6J18}\label{eq:6J18}
\end{equation*}
